%=====================================================================
\chapter{Security, NoSQL, and the Semester Project}\label{ch:project}
%=====================================================================
\locator{6}{Part VI --- Making It Work, and What Comes Next}

\begin{outcomes}
By the end of this chapter you will be able to:
\begin{tight}
  \item \textbf{Explain} SQL injection, \textbf{demonstrate} it against your
        own database, and \textbf{prevent} it with parameters.
  \item \textbf{Write} the \sqlkw{GRANT} and \sqlkw{REVOKE} statements a server
        would need for the university's external schemas, and \textbf{say} what
        SQLite offers instead.
  \item \textbf{Describe} the four families of NoSQL database and
        \textbf{state} what each gives up to get what it offers.
  \item \textbf{Store and query} a document inside SQLite using its JSON
        functions, and \textbf{compare} it with the relational form.
  \item \textbf{Decide} whether a given system needs a server DBMS, using the
        limits measured in this course.
\end{tight}
\end{outcomes}

\begin{prereq}
Appendix~A, section A.6 --- \textbf{the scope statement, which governs this
chapter's laboratory}. \chref{whodoeswhat} for the authorisation manager SQLite
lacks; \chref{views} for external schemas; \chref{transactions} for the
one-writer limit. The finished \texttt{university.db}.
\end{prereq}

\begin{keyterms}
SQL injection $\bullet$ parameterised query $\bullet$ placeholder $\bullet$
least privilege $\bullet$ \sqlkw{GRANT} $\bullet$ \sqlkw{REVOKE} $\bullet$
role $\bullet$ encryption at rest $\bullet$ NoSQL $\bullet$ key--value store
$\bullet$ document store $\bullet$ wide-column store $\bullet$ graph database
$\bullet$ schema-on-read $\bullet$ horizontal scaling $\bullet$ eventual
consistency $\bullet$ CAP $\bullet$ JSON
\end{keyterms}

\section{Security: the Component SQLite Does Not Have}

\chref{whodoeswhat} drew six components and said SQLite was missing one: the
authorisation manager. HEC's 2017 and 2023 outlines do not mention security at
all; HEC 2025's fifth learning outcome requires it, which is why it is here.
It has two halves --- keeping out people who should not be in, and making sure
that people who \emph{are} in cannot make the database do something it should
not --- and SQLite can help with only the second.

\begin{alertbox}[Scope --- this governs the laboratory below]
Everything in this chapter runs against \textbf{a database file you created,
on a machine you own}: your \texttt{university.db}, with invented data, which
you can delete. \textbf{Running any of it against a system you do not own and
have not been authorised in writing to test is a criminal offence in Pakistan}
under the \textbf{Prevention of Electronic Crimes Act 2016}, sections~3 and~4.
The purpose of the demonstration is that you never write the vulnerable code,
not that you can find it in somebody else's.
\end{alertbox}

\begin{definitionbox}[SQL injection]
\term{SQL injection} happens when a program builds a SQL statement by pasting
in text that a user supplied. If the text contains SQL, the database runs it
--- as part of the statement the programmer wrote, with the programmer's
permissions.

\smallskip
A \term{parameterised query} sends the statement and the values separately.
The statement has \term{placeholders} (\texttt{?}); the values are bound to
them afterwards and are \emph{never} parsed as SQL, whatever they contain.
\end{definitionbox}

\begin{worked}[One login check, built two ways]
A student portal checks a roll number before showing marks.

\medskip
\textbf{The vulnerable version} --- text pasted into SQL:
\begin{lstlisting}[language=python]
roll = input("Roll number: ")
sql = f"SELECT name FROM student WHERE roll_no = '{roll}'"
rows = cx.execute(sql).fetchall()
\end{lstlisting}

\textbf{Step 1 --- an ordinary input.} The roll number \texttt{2023-IT-00003}
produces the statement
\begin{lstlisting}[language=sql]
SELECT name FROM student WHERE roll_no = '2023-IT-00003'
\end{lstlisting}
and one row comes back. Correct.

\medskip
\textbf{Step 2 --- an input that is SQL.} \texttt{x' OR '1'='1} gives
\begin{lstlisting}[language=sql]
SELECT name FROM student WHERE roll_no = 'x' OR '1'='1'
\end{lstlisting}
The quote in the input closed the string the programmer opened; \texttt{OR
'1'='1'} is true for every row. On the seeded database it returns \textbf{all
4{,}000 students}.

\medskip
\textbf{Step 3 --- the parameterised version.}
\begin{lstlisting}[language=python]
rows = cx.execute("SELECT name FROM student WHERE roll_no = ?",
                  (roll,)).fetchall()
\end{lstlisting}
The same input now looks for a student whose roll number is, literally,
\texttt{x' OR '1'='1}. There is none: \textbf{0 rows}. The quote is a character
in a value, not the end of a string, because the value was never parsed.

\medskip
\textbf{The rule.} Never build SQL from input with \texttt{+}, \texttt{\%} or
an f-string. Every value goes through a placeholder. There is no escaping
function that is as safe, and no input validation that replaces it.
\end{worked}

\begin{examplebox}[The external schemas of Chapter 2, as a server would enforce them]
SQLite cannot run any of this. It is written out because it is what the
external level of \chref{architecture} \emph{means} once a database has users:

\begin{lstlisting}[language=sql]
CREATE ROLE fees_office;
CREATE ROLE exams_office;
CREATE ROLE portal;

GRANT SELECT ON fees_view TO fees_office;          -- names and rolls only
GRANT SELECT, UPDATE (grade) ON enrolment TO exams_office;
GRANT SELECT ON transcript TO portal;              -- a view, never a table

REVOKE ALL ON student FROM portal;                 -- least privilege
\end{lstlisting}

\smallskip
\term{Least privilege}: each user gets exactly what its job needs and no more.
Then even a successful SQL injection through the portal can only read what
\texttt{portal} may read --- the damage is bounded by the grant, not by the
programmer's care.

\smallskip
\textbf{What SQLite offers instead} is the operating system: whoever can read
the file can read all of it. File permissions, an encrypted disk, or an
encrypted SQLite build (SQLCipher) protect the file; nothing protects one
table from another.
\end{examplebox}

\section{NoSQL: What It Gives Up to Get What It Offers}

\begin{definitionbox}[Four families]
\term{NoSQL} names databases that do not use the relational model. The four
common families:

\smallskip
\begin{tabular}{@{}L{3.0cm}L{4.6cm}L{6.2cm}@{}}
\textbf{Family} & \textbf{Stores} & \textbf{Good at} \\
key--value & a value per key, opaque & caching, sessions: get and put by key \\
document & a nested document (JSON) per key & records read whole, with
  varying shape \\
wide-column & rows with many sparse columns & huge write volumes across many
  machines \\
graph & nodes and edges & paths: friends of friends, prerequisites of
  prerequisites \\
\end{tabular}

\smallskip
What they usually give up, to different degrees: \textbf{joins},
\textbf{a declared schema} (the data's shape is checked when read, not when
written --- \term{schema-on-read}), and \textbf{multi-record ACID
transactions}. What they gain is the ability to spread one database across
many machines (\term{horizontal scaling}).
\end{definitionbox}

\begin{worked}[One student's record as a document, inside SQLite]
SQLite has JSON functions, so a document store can be shown without
installing one.

\medskip
\textbf{Step 1 --- build the document.} One student, with their enrolments
nested inside:
\begin{lstlisting}[language=sql]
CREATE TABLE student_doc (roll_no TEXT PRIMARY KEY, doc TEXT NOT NULL);
INSERT INTO student_doc
SELECT s.roll_no,
       json_object('name', s.name, 'programme', s.programme,
         'enrolments', (SELECT json_group_array(
             json_object('course', e.course_code, 'term', e.semester,
                         'grade', e.grade))
            FROM enrolment e WHERE e.roll_no = s.roll_no))
  FROM student s;
\end{lstlisting}

\textbf{Step 2 --- what it is good at.} One student's whole record is one row,
read in one lookup, with no join:
\begin{lstlisting}[language=sql]
SELECT doc FROM student_doc WHERE roll_no = '2023-IT-00003';
\end{lstlisting}

\textbf{Step 3 --- what it is bad at.} ``How many A grades were given in
CS-101?'' was one \sqlkw{GROUP BY} in \chref{grouping}. Now every document must
be opened and unpacked:
\begin{lstlisting}[language=sql]
SELECT COUNT(*) FROM student_doc, json_each(doc, '$.enrolments') j
 WHERE json_extract(j.value, '$.course') = 'CS-101'
   AND json_extract(j.value, '$.grade')  = 'A';
\end{lstlisting}
On the seeded database both forms find \textbf{161}. The relational query took
9.8\,ms; the document form took \textbf{78.4\,ms}, eight times as long, because
it parsed every one of 4{,}000 documents to look inside each.

\textbf{Step 4 --- what it has given up.} Nothing stops a document with
\texttt{"grade": "Z"}, or \texttt{"grde"}, or a course that does not exist ---
\chref{constraints}'s three model-based constraints are gone, and a course's
title, if stored, is copied into every document that mentions it:
\chref{normalforms}'s update anomaly, by design.

\medskip
\textbf{The lesson.} A document store is the right model when records are read
and written \emph{whole} and rarely questioned \emph{across}. A university's
results are the opposite --- nearly every question in Part~IV crosses students
--- which is why this course is relational. That is a decision about the
\emph{questions}, not about which technology is newer.
\end{worked}

\dsfig{%
\begin{tikzpicture}[
  tb/.style={rounded corners=2pt, draw=#1, line width=0.9pt, fill=#1!7,
             text=#1!60!black, font=\scriptsize\ttfamily, align=left,
             inner sep=3pt},
  ttl/.style={font=\scriptsize\bfseries, text=#1, anchor=south},
  lk/.style={-{Stealth[length=2mm]}, neutral!60, line width=0.7pt},
  nt/.style={font=\scriptsize\itshape, text=#1, align=center}
]
  \node[ttl=primary] at (-4.0,2.2) {relational: the fact lives once};
  \node[tb=primary] (s) at (-5.3,1.3) {student\\2023-IT-00003\\Ayesha Q.};
  \node[tb=accent]  (e) at (-2.7,1.3) {enrolment\\00003 CS-101 A\\00003 MA-150 B\\\dots};
  \node[tb=greenhl] (c) at (-4.0,-0.35) {course\\CS-101 \ PH Course 1\\MA-150 \ \dots};
  \draw[lk] (e.west) -- (s.east);
  \draw[lk] (e.south) -- (c.north east);
  \node[nt=primary, text width=5.2cm] at (-4.0,-1.45)
       {``A grades in CS-101'': one index search on one table. 9.8\,ms};

  \node[ttl=purplehl] at (3.6,2.2) {document: the student lives once};
  \node[tb=purplehl, text width=5.4cm] (d) at (3.6,0.75)
       {\textbraceleft{} "name": "Ayesha Q.",\\
        \ \ "enrolments": [\\
        \ \ \ \ \textbraceleft"course":"CS-101","grade":"A"\textbraceright,\\
        \ \ \ \ \textbraceleft"course":"MA-150","grade":"B"\textbraceright,
        \dots ] \textbraceright};
  \node[nt=purplehl, text width=5.6cm] at (3.6,-0.85)
       {one lookup returns the whole record ---\\
        but ``A grades in CS-101'' opens all 4{,}000. 78.4\,ms};
  \node[nt=accent, text width=5.6cm] at (3.6,-1.6)
       {and nothing stops \texttt{"grde": "Z"}};
\end{tikzpicture}}%
{One student, held two ways. The relational form stores each fact once and
answers questions across students with a search; the document form stores each
student once and answers questions about one student in a single read. Which
is right depends on which kind of question the system is asked.}{docrel}

\begin{conceptbox}[When to leave SQLite --- using what this course measured]
\begin{tabular}{@{}L{4.4cm}L{4.4cm}L{5.0cm}@{}}
\textbf{Limit} & \textbf{Measured in} & \textbf{Move to a server when} \\
one writer at a time & \chref{transactions} & many users write at the same
  moment, all day \\
no users or privileges & \chref{whodoeswhat}, this chapter & different people
  must see different parts \\
on one machine & \chref{firstdatabase} & the data must be reached over a
  network by many programs \\
optimiser without STAT4 & \chref{indexes} & large, skewed tables where plan
  choice matters \\
\end{tabular}

\smallskip
And stay with it when none of those apply: SQLite is the most widely deployed
database in the world --- in every phone, every browser --- precisely because
for one program and one file it has nothing to administer. \chref{firstdatabase}
counted that as zero installation steps, zero ports and zero accounts; this
table is the other side of the same trade.
\end{conceptbox}

\section{The Semester Project}

\begin{definitionbox}[What you submit]
The project is the \texttt{university.db} you have built since
\chref{firstdatabase}, finished. Submit one folder containing:

\begin{tightnum}
  \item \texttt{design.pdf} --- your ER diagram with every cardinality and
        participation constraint (\chref{cardinality}), and the dependencies
        of every table (\chref{fds}).
  \item \texttt{v1\_create.sql}, \texttt{v2\_constraints.sql} and
        \texttt{v3\_final.sql} --- the schema and each migration
        (\chref{evolution}), each runnable from an empty file.
  \item \texttt{queries.sql} --- twelve queries answering the questions on
        the next page, each preceded by a comment stating the question.
  \item \texttt{triggers.sql} --- the six-section rule for \sqlkw{INSERT}
        \emph{and} \sqlkw{UPDATE} (\chref{views}), and an audit trigger.
  \item \texttt{report.pdf} --- six pages: three decisions you made and the
        evidence for each, including at least one measurement of your own.
  \item \texttt{university.db} itself, after \texttt{PRAGMA
        integrity\_check} and \texttt{PRAGMA foreign\_key\_check} both pass.
\end{tightnum}
\end{definitionbox}

\begin{examplebox}[The twelve questions]
\begin{tightnum}
  \item Every student's transcript, with course titles and credit hours.
  \item Each student's grade-point average over graded courses only.
  \item The top three students in each programme, ties included.
  \item Sections with no instructor --- the 73 of \chref{joins}.
  \item Instructors teaching nothing, written with \sqlkw{NOT EXISTS}.
  \item Students who have never received an A.
  \item Every prerequisite of a course, to any depth.
  \item Students enrolled in a course without having passed its direct
        prerequisites.
  \item Enrolments and average credit hours per department.
  \item Each course's pass rate, ungraded enrolments excluded and counted.
  \item A running total of credit hours per student, in correct term order.
  \item One question of your own that needs a window function.
\end{tightnum}

\smallskip
For each one, \texttt{queries.sql} must also contain its
\sqlkw{EXPLAIN QUERY PLAN}, and your report must name any query whose plan you
changed with an index, with both timings.
\end{examplebox}

\begin{conceptbox}[How it is marked]
\begin{tabular}{@{}L{6.0cm}R{1.6cm}L{6.2cm}@{}}
\textbf{Component} & \textbf{\%} & \textbf{What earns the marks} \\
design: ER diagram and dependencies & 20 & every constraint argued, not just
  drawn \\
schema and migrations & 20 & runs from empty; every constraint enforced,
  checked \\
the twelve queries & 30 & correct with nulls; \sqlkw{NOT IN} on a nullable
  column scores 0 \\
triggers & 10 & every event that can break the rule is covered \\
report and measurement & 20 & claims backed by numbers you produced \\
\end{tabular}

\smallskip
A query that returns the right answer on the seeded data and the wrong one when
a null appears is wrong. Test every one of them with a null.
\end{conceptbox}

\begin{pitfall}
\begin{tight}
  \item \textbf{Building SQL from input.} Parameters, always.
  \item \textbf{Escaping instead of parameterising.} Escaping is a list of
        characters somebody remembered; a parameter is never parsed.
  \item \textbf{Treating a view as access control in SQLite.} No users, so no
        restriction.
  \item \textbf{Choosing NoSQL because it is newer.} Choose by the questions:
        across records, relational; whole records, document.
  \item \textbf{Submitting a schema that only runs on your machine.} Run every
        script from an empty file before handing in.
  \item \textbf{Testing queries only on data without nulls.} The FALL2026
        grades are null for a reason.
\end{tight}
\end{pitfall}

\begin{lab}[Inject your own database, then make it impossible]
Read Appendix~A, section A.6, first. Every step uses \texttt{university.db}
on your own machine.

\begin{lstlisting}[language=bash]
# --- 1. The seeded database, which you can delete -------------------
python code/schema/seed.py university.db
\end{lstlisting}
\begin{lstlisting}[language=python]
# --- 2. The vulnerable query, in a file called inject.py ------------
import sqlite3
cx = sqlite3.connect("university.db")
for roll in ["2023-IT-00003", "x' OR '1'='1"]:
    sql = f"SELECT COUNT(*) FROM student WHERE roll_no = '{roll}'"
    print(repr(roll), "->", cx.execute(sql).fetchone()[0], "rows")

# --- 3. The same two inputs, parameterised --------------------------
for roll in ["2023-IT-00003", "x' OR '1'='1"]:
    n = cx.execute("SELECT COUNT(*) FROM student WHERE roll_no = ?",
                   (roll,)).fetchone()[0]
    print(repr(roll), "->", n, "rows (parameterised)")
\end{lstlisting}
\begin{lstlisting}[language=sql]
-- --- 4. A document, and a question it is bad at -----------------
.open university.db
CREATE TABLE student_doc (roll_no TEXT PRIMARY KEY, doc TEXT NOT NULL);
-- Type in Step 1 of the worked example, then:
SELECT json_extract(doc, '$.name'),
       json_array_length(doc, '$.enrolments')
  FROM student_doc LIMIT 3;

-- --- 5. Compare one cross-record question both ways --------------
.timer on
SELECT COUNT(*) FROM enrolment
 WHERE course_code = 'CS-101' AND grade = 'A';
SELECT COUNT(*) FROM student_doc, json_each(doc, '$.enrolments') j
 WHERE json_extract(j.value, '$.course') = 'CS-101'
   AND json_extract(j.value, '$.grade') = 'A';
-- Same count? Record both timings.

-- --- 6. Put in a document no relational table would accept --------
INSERT INTO student_doc VALUES ('X-1',
  '{"name":"Nobody","enrolments":[{"course":"NONE","grde":"Z"}]}');
-- Accepted. List the constraints from Chapter 8 it breaks.
\end{lstlisting}
\begin{lstlisting}[language=bash]
# --- 7. Run your project's checks before submitting ------------------
sqlite3 university.db "PRAGMA integrity_check; PRAGMA foreign_key_check;"
\end{lstlisting}

\smallskip
\textbf{What to notice.} Step 2's second line and step 3's second line are the
same input, giving 4{,}000 rows and 0. Nothing about the input changed; only
whether it was ever parsed as SQL. And step 6 is the other half of the
chapter: a model that accepts anything checks nothing.
\end{lab}

\begin{checkpoint}
\begin{tightnum}
  \item Explain exactly why \texttt{x' OR '1'='1} returned every student when
        pasted into the query and none when bound as a parameter, naming the
        character that made the difference and what happened to it in each
        case.
  \item Write the \sqlkw{GRANT} statements giving the examinations office
        permission to read every enrolment and change only grades, and say
        what SQLite would let that office do with the same file.
  \item For the question ``all of one student's grades'' and the question
        ``the number of A grades in CS-101'', say which model --- relational or
        document --- answers each more naturally, and why.
\end{tightnum}
\end{checkpoint}

\begin{chaptersummary}
\begin{tight}
  \item \textbf{SQL injection}: text pasted into SQL is parsed as SQL.
        \texttt{x' OR '1'='1} returned \textbf{4{,}000 students} pasted and
        \textbf{0} as a parameter.
  \item A \textbf{parameterised query} sends values separately; they are never
        parsed. There is no escaping function as safe.
  \item A server enforces external schemas with roles, \sqlkw{GRANT} and
        \sqlkw{REVOKE}, and \textbf{least privilege} bounds the damage of any
        flaw. SQLite has no users: the file's permissions are the only
        control.
  \item \textbf{NoSQL} families --- key--value, document, wide-column, graph
        --- usually give up joins, a declared schema and multi-record ACID to
        gain horizontal scaling.
  \item A document is read whole in one lookup and questioned across records
        only by unpacking every one; it accepts any shape, so none of
        \chref{constraints}'s constraints hold.
  \item Choose the model by the questions. A university's questions cross
        records, which is why this course is relational.
  \item Leave SQLite when the measured limits bite: many simultaneous writers,
        per-user access, network access, large skewed tables. Otherwise stay.
  \item The project is the \texttt{university.db} built all term, with its
        design, migrations, twelve queries, triggers and a report backed by
        your own measurements.
\end{tight}
\end{chaptersummary}

\begin{reviewq}
  \item \textbf{(MCQ)} A parameterised query prevents SQL injection because
        the value is: (a)~escaped; (b)~validated; (c)~never parsed as SQL;
        (d)~encrypted.
  \item \textbf{(MCQ)} Which NoSQL family suits ``friends of friends''?
        (a)~key--value; (b)~document; (c)~wide-column; (d)~graph.
  \item \textbf{(Short)} State the principle of least privilege and say how it
        limits the damage of an injection that does succeed.
  \item \textbf{(Short)} What does ``schema-on-read'' mean, and what does a
        document store therefore not check on write?
  \item \textbf{(Short)} Give two limits measured in this course that would
        justify moving the university's database to a server.
  \item \textbf{(Applied)} Rewrite a program of your own from the term that
        builds any SQL with string formatting so that every value is a
        parameter. Show one input that changed the old version's behaviour
        and does not change the new one's.
  \item \textbf{(Applied)} A startup proposes storing the university's
        records as one JSON document per student ``because it is simpler''.
        Using this chapter's worked example and the twelve project questions,
        say which questions become harder and why, what integrity is lost, and
        under what changed requirements the startup would be right.
\end{reviewq}
